% !TEX root = ../main.tex

\chapter{Background, Stato dell’Arte e Fondamenti}
\label{chap:background}

Questo capitolo fornisce il contesto teorico e tecnico necessario a inquadrare il lavoro presentato nella tesi.
In particolare, vengono introdotti il dominio applicativo delle \emph{New Psychoactive Substances} (NPS), le tecniche analitiche basate su spettrometria di massa a ionizzazione elettronica (EI), le principali fonti di dati spettrali utilizzate in ambito forense, nonché i modelli fisici e metodologici alla base del sistema PENTION-M.
Il capitolo si conclude con una panoramica sullo stato dell’arte relativo al \emph{Physics-Informed Machine Learning} e al ruolo dell’MLOps in sistemi cyber-fisici e contesti forensi.

\section{New Psychoactive Substances (NPS)}
\label{sec:nps}

Le \emph{New Psychoactive Substances} (NPS) rappresentano una classe eterogenea di composti chimici progettati per imitare gli effetti di sostanze stupefacenti tradizionali, eludendo al contempo i controlli normativi e le liste delle sostanze proibite.
Secondo le definizioni fornite dagli organismi internazionali di riferimento, quali EMCDDA ed UNODC, una NPS è una sostanza non controllata o recentemente controllata che può rappresentare un rischio per la salute pubblica \cite{emcdda_nps, unodc_nps}.

Dal punto di vista operativo, le NPS pongono sfide significative alle forze dell’ordine (LEAs) e ai laboratori forensi:
\begin{itemize}
    \item elevata variabilità strutturale e molecolare;
    \item rapida introduzione sul mercato di nuove varianti;
    \item somiglianza spettrale tra composti appartenenti alla stessa classe chimica;
    \item necessità di identificazione rapida in contesti non controllati.
\end{itemize}

Queste caratteristiche rendono particolarmente complessa l’identificazione affidabile delle NPS in scenari operativi dinamici, come controlli su strada, ispezioni portuali o interventi in aree urbane densamente popolate.

\section{Spettrometria di Massa a Ionizzazione Elettronica}
\label{sec:ei}

La spettrometria di massa a ionizzazione elettronica (Electron Ionization, EI) è una delle tecniche analitiche più consolidate in ambito forense per l’identificazione di sostanze chimiche.
Nel processo EI, le molecole vengono ionizzate mediante impatto elettronico ad alta energia (tipicamente 70 eV), generando una serie di frammenti caratteristici che producono uno spettro di massa discreto \cite{gross_mass_spec}.

Uno spettro EI può essere rappresentato come una funzione discreta:
\[
S = \{(m/z_i, I_i)\}_{i=1}^{N}
\]
dove $m/z_i$ indica il rapporto massa/carica e $I_i$ l’intensità relativa del frammento.

I principali vantaggi dell’EI in ambito NPS includono:
\begin{itemize}
    \item elevata riproducibilità degli spettri;
    \item compatibilità con ampie librerie di riferimento;
    \item robustezza rispetto a condizioni operative variabili.
\end{itemize}

Tuttavia, la frammentazione intensa introdotta dall’EI può generare spettri molto simili tra composti strutturalmente affini, rendendo necessarie tecniche avanzate di analisi e classificazione automatica.

\section{Dataset e Banche Dati per la Caratterizzazione delle NPS}
\label{sec:datasets}

L’efficacia di un sistema di classificazione basato su spettrometria di massa dipende in larga misura dalla qualità, dalla copertura e dalla coerenza dei dataset utilizzati per l’addestramento dei modelli.

Nel contesto delle NPS, non esiste una singola banca dati esaustiva: i dati disponibili sono distribuiti su più fonti, spesso eterogenee per formato, qualità e convenzioni di naming.
Questa frammentazione rende necessaria una fase esplicita di integrazione e normalizzazione dei dataset.

\subsection{Panoramica delle Fonti di Dati EI}
\label{subsec:ei_sources}

Le principali fonti di dati EI utilizzate in questo lavoro includono:
\begin{itemize}
    \item librerie open access (es. SWGDRUG, MoNA);
    \item database accademici e istituzionali (es. SDBS);
    \item dataset rilasciati come supplementi a pubblicazioni scientifiche;
    \item archivi forensi ad accesso regolamentato (es. ENFSI).
\end{itemize}

L’uso combinato di tali fonti consente di ottenere una copertura ampia delle classi NPS, a fronte però di una maggiore complessità nella fase di preprocessing.

\subsection{SWGDRUG}
\label{subsec:swgdrug}

La libreria SWGDRUG (\emph{Scientific Working Group for the Analysis of Seized Drugs}) rappresenta uno degli standard de facto per la spettrometria di massa forense.
Essa fornisce spettri EI di elevata qualità, generalmente accompagnati da metadati affidabili e coerenti \cite{swgdrug}.

I principali limiti di SWGDRUG riguardano:
\begin{itemize}
    \item la copertura incompleta delle NPS più recenti;
    \item la presenza di varianti nominali per lo stesso composto.
\end{itemize}

\subsection{MoNA -- MassBank of North America}
\label{subsec:mona}

MoNA è una piattaforma open access che aggrega spettri di massa provenienti da numerose fonti \cite{mona}.
Pur offrendo un’elevata quantità di dati, MoNA presenta una maggiore eterogeneità nella qualità degli spettri e nella completezza dei metadati, rendendo necessaria una selezione mirata dei record EI.

\subsection{SDBS}
\label{subsec:sdbs}

Il database SDBS (\emph{Spectral Database for Organic Compounds}) fornisce spettri di riferimento per numerosi composti organici \cite{sdbs}.
Nel contesto di questa tesi, SDBS è stato utilizzato per il recupero mirato di spettri EI relativi a specifiche sostanze NPS, in particolare nei casi di assenza di dati in altre librerie.

\subsection{Problemi di Merging e Normalizzazione}
\label{subsec:merging}

L’integrazione di più fonti EI richiede la risoluzione di diversi problemi:
\begin{itemize}
    \item duplicazione di composti con naming differente;
    \item differenze nella risoluzione $m/z$;
    \item normalizzazione delle intensità;
    \item gestione di spettri incompleti o rumorosi.
\end{itemize}

Nel sistema PENTION-M, gli spettri EI sono stati vettorizzati su un dominio discreto $m/z \in [1,600]$, con normalizzazione delle intensità e deduplicazione basata su similarità semantica dei nomi chimici e, ove disponibile, su identificatori strutturali.

\subsection{Dataset Finale: \texttt{PENTION\_EI\_Complete}}
\label{subsec:pention_dataset}

Il risultato del processo di integrazione è il dataset \texttt{PENTION\_EI\_Complete}, che rappresenta una base spettrale unificata per l’addestramento del classificatore NPS.
Il dataset copre un ampio insieme di classi chimiche rilevanti dal punto di vista forense ed è progettato per supportare modelli di classificazione robusti e generalizzabili.

\section{Simulazione della Dispersione Chimica in Ambienti Urbani}
\label{sec:dispersion}

La modellazione della dispersione di sostanze chimiche in atmosfera è un problema classico della fisica ambientale.
In contesti urbani, tale problema è ulteriormente complicato dalla presenza di edifici, dalla variabilità meteorologica e dalla complessità geometrica dello spazio.

I modelli gaussiani di dispersione, in particolare i modelli \emph{Gaussian Puff} e \emph{Gaussian Plume}, rappresentano una soluzione consolidata per la simulazione di rilasci puntuali \cite{turner_gaussian}.
Nel lavoro presentato, tali modelli sono utilizzati come base fisica per la generazione di mappe di concentrazione coerenti, successivamente corrette mediante tecniche di apprendimento automatico informate dalla fisica.

\section{Physics-Informed Machine Learning}
\label{sec:piml}

Il \emph{Physics-Informed Machine Learning} (PIML) rappresenta un paradigma che integra conoscenza fisica nei modelli di apprendimento automatico al fine di migliorarne robustezza, interpretabilità e capacità di generalizzazione \cite{kashinath_piml, meng_piml}.

Nel PIML, la conoscenza fisica può essere incorporata a diversi livelli:
\begin{itemize}
    \item livello dei dati (dataset generati da modelli fisici);
    \item livello delle feature (variabili fisicamente significative);
    \item livello della funzione di perdita (vincoli fisici).
\end{itemize}

Nel contesto di PENTION-M, il PIML è utilizzato per correggere la dispersione gaussiana e per migliorare la localizzazione della sorgente, garantendo coerenza con i principi fisici del fenomeno simulato.

\section{MLOps in Sistemi Cyber-Fisici e Contesti Forensi}
\label{sec:mlops}

L’adozione di modelli di Machine Learning in sistemi operativi richiede pipeline MLOps in grado di gestire versionamento, monitoraggio, drift dei dati e auditabilità \cite{faubel_mlops, calefato_secmlops}.

In contesti forensi e investigativi, tali requisiti sono ulteriormente rafforzati dalla necessità di:
\begin{itemize}
    \item tracciabilità end-to-end delle decisioni;
    \item riproducibilità delle inferenze;
    \item garanzia di integrità e non ripudio dei risultati.
\end{itemize}

Il sistema PENTION-M integra tali principi mediante una pipeline MLOps containerizzata e un meccanismo di logging forense basato su firme digitali e bundle verificabili.

\section{PENTION-S: Architettura di Riferimento}
\label{sec:pention_s}

PENTION-S rappresenta la prima iterazione del progetto PENTION, focalizzata su un’architettura statica di analisi e classificazione delle NPS.
Pur fornendo una base solida, PENTION-S presenta limiti intrinseci in termini di mobilità, adattabilità e integrazione operativa.

Questi limiti costituiscono la motivazione principale per la progettazione di PENTION-M, che estende il sistema verso un paradigma mobile, fisicamente informato e forense-ready, come discusso nei capitoli successivi.
